\section{About the Committee}

\subsection{Mandate}

As per the Constitution, the U.S. House of Representatives makes and passes federal laws. The House is one of Congress’s two chambers (the other is the U.S. Senate) and part of the federal government’s legislative branch. The number of voting representatives in the House is fixed by law at no more than 435, with representation proportional to the populations of the 50 states.\footnote{U.S. House of Representatives. (n.d.). \emph{The House explained}. House.gov. \url{https://www.house.gov/the-house-explained}}

\subsubsection{Primary Responsibilities}

\begin{itemize}
  \item \textbf{Legislation:} The House’s exclusive powers include proposing, debating, and passing bills.
  \item \textbf{Revenue bills:} Originating all bills regarding taxes and raising revenue.
  \item \textbf{Impeachment:} Holding the sole power to impeach federal officers, including the president.
  \item \textbf{Electoral contingency:} Electing the president if no candidate receives a majority of votes in the Electoral College.
\end{itemize}

\subsection{History}

The House of Representatives was established by Article One of the United States Constitution and first convened in 1789, initially meeting in New York City before the federal government moved to Philadelphia and ultimately to Washington, D.C.\footnote{U.S. House of Representatives, Office of the Historian. (n.d.). \emph{History of the House}. House.gov. \url{https://www.house.gov/the-house-explained/history-of-the-house}} It was designed as the chamber of Congress most directly accountable to the people, with members chosen by popular vote every two years, in contrast to the Senate (whose members were originally chosen by state legislatures until the Seventeenth Amendment in 1913 provided for their direct election).\footnote{American Academy of Arts and Sciences. (n.d.). \emph{The case for enlarging the House of Representatives}. \url{https://www.amacad.org/ourcommonpurpose/enlarging-the-house/section/2}}

As the nation expanded westward and its population grew, the size of the House increased after successive censuses.\footnote{U.S. House of Representatives, Office of the Historian. (n.d.). \emph{Determining apportionment}. History, Art \& Archives, U.S. House of Representatives. \url{https://history.house.gov/Institution/Apportionment/Determining-Apportionment/}} The Reapportionment Act of 1929 (building on the Permanent Apportionment Act) ultimately fixed the number of voting members at 435, where it has remained, with seats reapportioned among the states after each decennial census.\footnote{U.S. House of Representatives, Office of the Historian. (n.d.). \emph{The Permanent Apportionment Act of 1929}. History, Art \& Archives, U.S. House of Representatives. \url{https://history.house.gov/Historical-Highlights/1901-1950/The-Permanent-Apportionment-Act-of-1929/}}

Over time, the franchise broadened considerably: property and other restrictions fell away, the Fifteenth, Nineteenth, Twenty-Fourth, and Twenty-Sixth Amendments and the Voting Rights Act of 1965 progressively extended voting rights across race, sex, and age.\footnote{American Academy of Arts and Sciences. (n.d.). \emph{The case for enlarging the House of Representatives}. \url{https://www.amacad.org/ourcommonpurpose/enlarging-the-house/section/2}} The House has developed an elaborate committee system and body of standing rules to manage an ever-growing legislative workload, and it has played a central historical role through its powers over revenue, impeachment, and the shaping of national policy.\footnote{U.S. House of Representatives, Office of the Historian. (n.d.). \emph{House committees} [Fact sheet]. History, Art \& Archives, U.S. House of Representatives. \url{https://history.house.gov/Education/Fact-Sheets/Committees-Fact-Sheet2/}}

\subsection{Representatives}

Each representative is elected for a two-year term, serving the people of a specific congressional district. Among other duties, they introduce bills and resolutions, offer amendments, and serve on committees. The number of representatives with full voting rights has been 435 since 1913.\footnote{U.S. House of Representatives. (n.d.). \emph{The House explained}. House.gov. \url{https://www.house.gov/the-house-explained}}

The House’s composition was established by Article One of the United States Constitution. The House is composed of representatives who, pursuant to the Uniform Congressional District Act, sit in single-member congressional districts allocated to each state based on population as measured by the United States census, provided that each state gets at least one representative.\footnote{Congressional Research Service. (n.d.). \emph{Election policy fundamentals: Single-member House districts} (IF12567). Congress.gov. \url{https://www.congress.gov/crs-product/IF12567}} The number of representatives per state is proportional to population. Special elections may be held in the event of a vacancy.

Since its inception in 1789, all representatives have been directly elected. Although suffrage was initially limited, it gradually widened, particularly after the ratification of the Nineteenth Amendment and the civil rights movement.\footnote{American Academy of Arts and Sciences. (n.d.). \emph{The case for enlarging the House of Representatives}. \url{https://www.amacad.org/ourcommonpurpose/enlarging-the-house/section/2}}

\subsection{Committees}

The House’s standing committees have different legislative jurisdictions. Each considers bills and recommends measures for consideration by the House.\footnote{U.S. House of Representatives. (n.d.). \emph{The House explained}. House.gov. \url{https://www.house.gov/the-house-explained}}

\begin{itemize}
  \item They oversee agencies, programmes, and activities within their jurisdictions.
  \item The Committee of the Whole House is one in which all the members serve.
  \item Before members are assigned to committees, each committee’s size and the proportion of Democrats to Republicans must be decided by the party leaders.
  \item The total number of committee slots allotted to each party is approximately the same as the ratio between the majority party and the minority party members in the full chamber.
\end{itemize}

\subsection{Commissions}

Congress has created a wide variety of temporary and permanent commissions to serve as advisory bodies for investigative or policy-related issues, or to carry out administrative, interparliamentary, or commemorative tasks. These are composed of House members, private citizens, or a mix of both. In some cases, the commissions are entities of the House or Congress itself.\footnote{U.S. House of Representatives. (n.d.). \emph{The House explained}. House.gov. \url{https://www.house.gov/the-house-explained}}

\textbf{Examples:} Financial Crisis Inquiry Commission (temporary); Commission on Security and Cooperation in Europe (Helsinki Commission); House Page Board (fixed).

\subsection{Special Rules of Procedure}

The House of Representatives is governed by standing rules that make certain matters and actions privileged for consideration. On a day-to-day basis, the House can also decide to grant individual bills privileged access to the floor. The standing rules of the House include parliamentary mechanisms to use on bills and resolutions. Which of these will be employed depends on the extent to which Members wish to debate and amend the legislation.

\subsubsection{Key Procedural Mechanisms}

\begin{enumerate}
  \item \textbf{Privileged Status and Floor Access.} The House has rules that control what gets debated. Its standing rules determine which matters have priority (“privileged” status). The House can also give a specific bill special access to the floor using parliamentary procedures, allowing it to be considered sooner.\footnote{Congressional Research Service. (2019). \emph{The legislative process on the House floor: An introduction} (95-563). Congress.gov. \url{https://www.congress.gov/crs-product/95-563}}
  \item \textbf{Methods of Consideration.} The House has several ways to consider legislation, depending on how much debate is desired:
\end{enumerate}

A.   \textbf{Suspension of the Rules} (most common method; often used for non-controversial bills): Debate is limited; no amendments can be offered from the floor; requires a two-thirds majority of those voting to pass.\footnote{Ibid.}

B.    \textbf{Unanimous Consent:} Every Member agrees to consider the measure. Usually used for routine or widely supported matters. Any single Member can object and block this procedure.\footnote{Congressional Research Service. (n.d.). \emph{How measures are brought to the House floor: A brief introduction} (RS20067). Congress.gov. \url{https://www.congress.gov/crs-product/RS20067}}

C.    \textbf{Rules Committee Process:} The House Rules Committee decides how major bills will be debated. This gives House leaders significant control over the legislative process. The House may allow many amendments, allow only certain amendments, or prohibit amendments entirely.\footnote{Congressional Research Service. (2025). \emph{Special rules in the House of Representatives: Purpose and content} (R48308). Congress.gov. \url{https://www.congress.gov/crs-product/R48308}}

\begin{enumerate}
  \item \textbf{Committee of the Whole.} Many important bills are first debated in the Committee of the Whole, essentially the entire House operating under more flexible rules. Advantages include easier debate procedures, more opportunities for amendments, and faster consideration of legislation. After debate, the bill returns to the full House for final passage by a simple majority.\footnote{Congressional Research Service. (n.d.). \emph{How measures are brought to the House floor: A brief introduction} (RS20067). Congress.gov. \url{https://www.congress.gov/crs-product/RS20067}}
  \item \textbf{Quorum Requirements.} A quorum is the minimum number of Members required to conduct business. To allow business to proceed efficiently, a quorum is generally presumed to exist; Members do not constantly check whether enough Representatives are present. A quorum issue is usually raised only during votes.\footnote{Congressional Research Service. (n.d.). \emph{Voting and quorum procedures in the House of Representatives} (98-988). Congress.gov. \url{https://www.congress.gov/crs-product/98-988}}
  \item \textbf{Voting Methods.} The House uses several voting methods: Voice Vote (Members verbally say “Aye” or “No”); Division Vote (Members physically indicate their vote, and the votes are counted); Recorded Electronic Vote (if enough Members request it, an official electronic vote is taken, and each Member’s vote is recorded, the most formal method).\footnote{Ibid.}
  \item \textbf{Bicameral Agreement.} The U.S. Constitution requires that both the House and the Senate approve identical text of a bill before it may be sent to the President. If the Senate changes a House bill (or vice versa), the differences must be resolved.\footnote{Congressional Research Service. (2019). \emph{The legislative process on the House floor: An introduction} (95-563). Congress.gov. \url{https://www.congress.gov/crs-product/95-563}}
  \item \textbf{Resolving Differences Between Chambers.} Two principal methods exist:
  \item \textbf{Exchange of Amendments (“Ping-Pong”):} The bill moves back and forth between the House and Senate; each chamber proposes changes until both agree.
  \item \textbf{Conference Committee:} A special committee with Members from both chambers negotiates a compromise version. Once an agreement is reached, the House votes on it, the Senate votes on it, and if both approve, the bill goes to the President.
\end{enumerate}

\emph{Note: The House can amend its rules unilaterally; it need not consult with either the Senate or the President. The House is free to suspend, waive, or ignore its rules whenever it chooses to do so. By and large, the Speaker or whatever Representative is presiding usually does not enforce the rules at his or her own initiative. Instead, Members must protect their own rights by affirmatively making points of order whenever they believe the rules are about to be violated. House rules include several formal procedures for waiving or suspending certain other rules, and almost any rule can be waived by unanimous consent. Thus, the requirements and restrictions generally apply only if the House chooses to enforce them.}
